\chapter{วิศวกรรมการประมวลผลสตรีมภาพและเซนเซอร์แบบเรียลไทม์}
\label{ch:realtime_stream_engineering}

\section{บทนำและโจทย์ความท้าทายของการประมวลผลบนอุปกรณ์พกพา}

การพัฒนาแอปพลิเคชันเพื่อการวิเคราะห์คุณสมบัติดินจากวิดีโอสตรีมสด (Live Video Stream) ร่วมกับสัญญาณเซนเซอร์แบบเรียลไทม์ จำเป็นต้องเผชิญกับข้อจำกัดด้านทรัพยากรการประมวลผลและหน่วยความจำของสมาร์ทโฟน หากกระบวนการแปลงรหัสสีและการสกัดคุณลักษณะกระทำบนเธรดหลักของส่วนต่อประสานผู้ใช้ (UI Main Thread) จะส่งผลให้เกิดการกระตุกของหน้าจอ (Frame Drop) และเกิดความล่าช้าสะสม บทนี้จึงนำเสนอการออกแบบเชิงวิศวกรรมซอฟต์แวร์ขั้นสูงเพื่อบรรลุสมรรถนะการทำงานแบบ Zero-Latency

\section{สถาปัตยกรรมท่อส่งสตรีมภาพและเทคนิคการลดความซ้ำซ้อน}

ในระบบปฏิบัติการแอนดรอยด์ สตรีมภาพสดจากฮาร์ดแวร์กล้องจะถูกส่งมอบผ่านอินเทอร์เฟซในรูปแบบฟอร์แมต $\text{YUV420\_888}$ ซึ่งประกอบด้วยระนาบความสว่าง ($Y$) และระนาบสีคู่ควบ ($U, V$) การแปลงระนาบภาพขนาดความละเอียดเต็มเข้าสู่ปริภูมิสี $\text{sRGB}$ ทุกพิกเซลในทุกเฟรมจะสิ้นเปลืองพลังงานประมวลผลอย่างมหาศาล ระบบจึงใช้เทคนิคการสุ่มตัวอย่างแบบก้าวกระโดดเชิงพื้นที่ (Spatial Subsampling Step $k=2$ หรือ $k=4$) เฉพาะภายในขอบเขต $\text{ROI}$ เท่านั้น

\begin{equation}
\Omega_{\text{subsample}} = \left\{ (x, y) \in \Omega_{\text{ROI}} \;\middle|\; x \equiv 0 \pmod k \text{ และ } y \equiv 0 \pmod k \right\}
\end{equation}

เทคนิคดังกล่าวช่วยลดจำนวนพิกเซลที่ต้องประมวลผลลงถึง $75\%$ ถึง $93.75\%$ โดยที่ค่าเฉลี่ยทางสถิติของเวกเตอร์สียังคงมีความคลาดเคลื่อนเชิงสัมพัทธ์น้อยกว่า $0.4\%$ เมื่อเทียบกับการคำนวณแบบเต็มพิกเซล

\section{การจัดการหน่วยความจำแบบ Zero-Copy และ Isolate Multi-threading}

เพื่อป้องกันปัญหาตัวเก็บขยะหน่วยความจำ (Garbage Collector Spike) ที่มักเกิดขึ้นจากการจองและปลดปล่อยอาร์เรย์พิกเซลซ้ำๆ ในภาษารันไทม์ ระบบจึงใช้พูลหน่วยความจำแบบหมุนเวียน (Circular Byte Buffer Pool) ที่สร้างไว้ล่วงหน้า (Pre-allocated Memory) ร่วมกับการถ่ายโอนงานคำนวณหนักไปยังเธรดแยกอิสระ (Background Isolate Worker) ผ่านกลไก $\text{TransferableTypedData}$

\begin{theorybox}
\textbf{วิศวกรรมการประมวลผลแบบอะซิงโครนัสขนานบนสถาปัตยกรรม Flutter}\par
กระบวนการรับเฟรมภาพสด ($60\text{ FPS}$) และการรับข้อมูลเซนเซอร์ผ่านบลูทูธพลังงานต่ำ ($\text{BLE } 10\text{ Hz}$) ดำเนินการแยกเธรดกันอย่างเด็ดขาด โดยข้อมูลภาพจะถูกประมวลผลในไอโซเลตของคอมพิวเตอร์วิทัศน์ ในขณะที่การประสานเวลาอาศัยหน้าต่างเวลาสัมพัทธ์ (Time-window Synchronization) ซึ่งเชื่อมโยงข้อมูลสองชุดที่บันทึกห่างกันไม่เกิน 100 มิลลิวินาทีเข้าเป็นหนึ่งตัวอย่างการอนุมาน
\end{theorybox}

\section{การประสานเวลาและการกรองสัญญาณรบกวนเชิงสถิติ}

เนื่องจากตัวอย่างสารละลายหรือหน้าดินอาจมีฟองอากาศ หรือการสั่นไหวของมือผู้ใช้งานขณะถือสมาร์ทโฟน ระบบจึงติดตั้งตัวกรองค่าเฉลี่ยเคลื่อนที่แบบถ่วงน้ำหนักตามเวลา (Exponential Moving Average หรือ EMA) สำหรับผลลัพธ์การทำนายในแต่ละเฟรม

\begin{equation}
\mathbf{y}_{\text{smooth}}^{(t)} = \beta \cdot \mathbf{y}_{\text{smooth}}^{(t-1)} + (1 - \beta) \cdot \widehat{\mathbf{y}}^{(t)}
\end{equation}

โดยที่ $\beta \in [0.7, 0.85]$ เป็นแฟกเตอร์การหน่วงเวลา ซึ่งช่วยให้เส้นกราฟและตัวเลขค่าที่แสดงบนหน้าปัดมีความนิ่งและอ่านง่าย สบายตา โดยที่ยังคงตอบสนองต่อการเปลี่ยนแปลงทางกายภาพจริงได้อย่างรวดเร็ว
